\section{Wprowadzenie}

Tekst, według którego ktoś działa, ma mieć 1 znaczenie. Polska gramatyka zostawia
czytelnikowi wybór w stałych miejscach: brak rodzajników, swobodny szyk, dopełniacz,
spójniki, granice przedziałów. Zdanie \pkq{Dzieci do 3 lat wchodzą bezpłatnie.} ma
2 odczytania: jedno obejmuje dziecko, które ma mniej niż 3 lata; drugie obejmuje dziecko,
które ma mniej niż 4 lata.

Polszczyzna Kontrolowana, w skrócie PK, zamyka te miejsca wyboru \cite{pk}. PK w wersji 0.1
to zwykła polszczyzna oraz \pkRules{} numerowanych reguł w 5 grupach. Standard ma także
moduł F dla raportów i listę słów kluczowych wymagań. Wzorami PK są angielski standard
ASD-STE100 \cite{ste}, słowa kluczowe RFC 2119 \cite{rfc2119}, notacja EARS \cite{ears}
i polskie \pkq{Zasady techniki prawodawczej} \cite{ztp}.

Czytelnikiem tekstu bywa model językowy. Model wykonuje instrukcję, streszcza
regulamin i sprawdza wymaganie. Praca bada tezę z tytułu: PK zmniejsza złożoność
obliczeniową pracy z tekstem. Praca rozdziela tę tezę na 4 pytania.
\begin{enumerate}[label=(\roman*),leftmargin=2.4em]
  \item Jak zapisać PK w języku logiki i teorii mnogości?
  \item Które koszty obliczeniowe PK zmniejsza, a które koszty PK zwiększa?
  \item Ile wynosi narzut długości PK i jak modele językowe odczytują zdania PK?
  \item Co PK mówi o hipotezie Sapira-Whorfa dla człowieka i dla modelu językowego?
\end{enumerate}

\paragraph{Odpowiedzi w skrócie.}
Język jest parą: zbiór zdań i relacja interpretacji. Reguła PK jest operatorem, który usuwa
zdania albo usuwa odczytania (sekcja \ref{sec:model}). PK zmniejsza koszt 3 mechanizmami:
zwija zbiór odczytań, ogranicza rozmiar zdania i daje zdaniu lokalność
(sekcja \ref{sec:complexity}). PK ma 1 koszt: tekst rośnie. Mediana narzutu wynosi
\pkTokMedian{} \% tokenów (sekcja \ref{sec:results}). Spójność zbioru wymagań bez przeczeń
jest rozstrzygalna w czasie liniowym. Reguła C5 dopuszcza 1 przeczenie w zdaniu, a 1
przeczenie wystarcza do NP-zupełności (twierdzenie \ref{thm:np}). Wynik pomiaru na
modelach językowych podaje sekcja \ref{sec:probe}. Refleksję filozoficzną zawiera
sekcja \ref{sec:philosophy}.

\paragraph{Status wyniku.}
Twierdzenia o parsowaniu, o kodowaniu i o klauzulach Horna są znane; praca przywołuje
źródła. Wkładem pracy są: model mnogościowy reguł PK, klasyfikacja \pkRules{} reguł,
twierdzenie graniczne o 1 przeczeniu, zbiór \pkItems{} pozycji testowych, sprawdzarka reguł
oraz 5 pomiarów. Pomiary mają ograniczony zasięg; sekcja \ref{sec:threats} wymienia
zagrożenia dla trafności. Praca opisuje PK 0.1 i nie ocenia innych standardów.

\paragraph{Konwencje tekstu.}
Autor napisał tę pracę w PK. Legenda: \pkq{lub} znaczy co najmniej 1 z członów;
\pkq{albo} znaczy dokładnie 1 z członów. Część dziesiętną liczby oddziela przecinek.
Cytat, przykład zdania zwykłego i tytuł stoją w cudzysłowie; taki fragment leży poza PK.
Każda liczba w tekście pochodzi ze skryptu \texttt{scripts/measure.py}. Sprawdzarka
z sekcji \ref{sec:check} skontrolowała tekst pracy; wynik podaje tabela \ref{tab:corpus}.

\section{Stan wiedzy}
\label{sec:related}

\paragraph{Języki kontrolowane.}
Język kontrolowany to podzbiór języka naturalnego z ograniczoną gramatyką i ograniczonym
słownictwem. Kuhn opisał 100 języków kontrolowanych opartych na angielskim i zaproponował
klasyfikację PENS \cite{kuhn2014}. Klasyfikacja PENS ocenia 4 cechy: precyzję,
ekspresywność, naturalność i prostotę. Kuhn stwierdza o językach klasy P$^1$, czyli
o językach naturalnych, że automatyczna interpretacja jest problemem \pkq{AI-zupełnym}.
Precyzja języka decyduje więc o koszcie interpretacji.

Języki kontrolowane mają 2 rodowody. Pierwszy rodowód to czytelność dla człowieka:
\pkq{Basic English} Ogdena \cite{ogden1930} i przewodnik AECMA z roku 1986, z którego
powstał standard ASD-STE100 \cite{ste}. Drugi rodowód to interpretacja przez maszynę:
Attempto Controlled English ma gramatykę formalną i tłumaczenie na logikę pierwszego rzędu
\cite{fuchs2008}. PK 0.1 należy do pierwszego rodowodu. Standard PK 0.1 składa się z reguł
zapisanych słowami i nie zawiera gramatyki formalnej. W klasyfikacji PENS autor umieszcza
PK 0.1 w klasie P$^2$: język ma mniej niejednoznaczności niż język naturalny i nie ma
podstawy formalnej.

\paragraph{Kontekst polski.}
Polska tradycja zna reguły redagowania tekstów prawnych \cite{ztp} i ruch prostego języka;
norma ISO 24495-1 opisuje prosty język \cite{iso24495}. Prosty język służy łatwemu
zrozumieniu, a PK służy jednemu znaczeniu. Ziembiński opisuje amfibolię, czyli wieloznaczność
składniową zdania \cite{ziembinski2002}. Przybyszewski zbadał zdania przeczące z kwantyfikatorem
ogólnym \cite{przybyszewski2020}. Kapłański opisał kontrolowaną angielszczyznę
jako interfejs baz wiedzy \cite{kaplanski2011}. Autor nie znalazł opublikowanego standardu
kontrolowanej polszczyzny ogólnego przeznaczenia. Przegląd objął wyszukiwarki naukowe
w listopadzie 2026 r. i ma ograniczony zasięg.

\paragraph{Złożoność a kształt języka.}
Church i Patil pokazali, że liczba rozbiorów zdania z ciągiem fraz przyimkowych rośnie jak
liczby Catalana \cite{church1982}. Earley podał algorytm parsowania gramatyk bezkontekstowych
w czasie rzędu $n^3$, a dla gramatyk jednoznacznych w czasie rzędu $n^2$ \cite{earley1970}.
Knuth zdefiniował gramatyki LR($k$), które parser czyta w czasie liniowym \cite{knuth1965};
Leo rozszerzył algorytm Earleya na czas liniowy dla każdej gramatyki LR($k$) \cite{leo1991}.
Barton wykazał, że parsowanie gramatyk ID/LP, które opisują swobodny szyk, jest NP-zupełne
\cite{barton1985}. Pratt-Hartmann i Third zbadali fragmenty języka angielskiego: problem
spełnialności przechodzi od czasu wielomianowego do nierozstrzygalności, gdy fragment
dostaje kolejne konstrukcje \cite{pratt2004,pratt2006}. Wyniki te pokazują wspólną zasadę:
zakres dopuszczonych konstrukcji wyznacza klasę złożoności.

\paragraph{Niejednoznaczność a modele językowe.}
Zbiór AmbigQA zawiera pytania z więcej niż 1 odpowiedzią \cite{min2020}. Liu i współautorzy
zbudowali zbiór AmbiEnt i stwierdzili, że rozpoznanie niejednoznaczności jest dla modeli
trudne \cite{liu2023}. Farquhar i współautorzy mierzą niepewność modelu entropią liczoną
na poziomie znaczeń \cite{farquhar2024}. Piantadosi, Tily i Gibson bronią
niejednoznaczności: kontekst rozstrzyga znaczenie, więc język zyskuje krótsze formy
\cite{piantadosi2012}. Wynik Piantadosiego wyznacza cenę PK: tekst bez oparcia w kontekście
jest dłuższy.

\section{Model logiczno-mnogościowy}
\label{sec:model}

\subsection{Język jako para}

\begin{definition}
Niech $W$ oznacza zbiór sytuacji. \emph{Znaczenie} to podzbiór zbioru $W$. Zbiór znaczeń to
$M=\mathcal{P}(W)$. \emph{Język} to para $(L,R)$: $L$ jest zbiorem zdań, a
$R\subseteq L\times M$ jest \emph{relacją interpretacji}.
\end{definition}

\begin{definition}
\emph{Zbiór odczytań} zdania $s$ to $\sem{s}=\{m\in M:(s,m)\in R\}$. \emph{Stopień
niejednoznaczności} zdania $s$ to liczba $a(s)=|\sem{s}|$. Zdanie $s$ jest
\emph{jednoznaczne}, gdy $a(s)=1$.
\end{definition}

Język $(L,R)$ jest jednoznaczny dokładnie wtedy, gdy relacja $R$ jest funkcją z $L$ do $M$.
Zwykła polszczyzna jest relacją, która nie jest funkcją. Rysunek \ref{fig:relation}
pokazuje 2 zdania zwykłe z 4 odczytaniami i 4 zdania PK, z których każde ma 1 odczytanie.

\rysunek{relation}{Teoria mnogości: relacja interpretacji zwykłej polszczyzny i funkcja,
do której dąży PK.}{fig:relation}

\subsection{Reguła jako operator}

Standard PK zawiera reguły 2 rodzajów. Reguła pierwszego rodzaju zabrania konstrukcji:
reguła A3 zabrania zaimka, reguła B5 zabrania imiesłowu. Reguła drugiego rodzaju zostawia
zdanie i ustala odczytanie. Według reguły B2 pierwszy rzeczownik jest podmiotem. Według
reguły C1 słowo \pkq{lub} znaczy co najmniej 1 człon.

\begin{definition}
\emph{Ograniczenie} to operator $\rho_\varphi(L,R)=(L\cap\varphi,\;R\cap(\varphi\times M))$
dla zbioru zdań $\varphi$. \emph{Konwencja} to operator $\kappa_K(L,R)=(L,\;R\cap K)$ dla
relacji $K\subseteq L\times M$.
\end{definition}

\begin{proposition}
\label{prop:kernel}
Każde ograniczenie i każda konwencja zmniejsza język, zachowuje zawieranie i jest
idempotentna. Każde 2 operatory są przemienne. Złożenie reguł $r_1,\dots,r_n$ daje język
\[
L'=L\cap\bigcap_i\varphi_i,\qquad R'=R\cap\bigcap_j K_j\cap(L'\times M).
\]
\end{proposition}

\begin{proof}
Oba operatory są przecięciem z ustalonym zbiorem. Przecięcie zbiorów jest przemienne,
łączne i idempotentne.
\end{proof}

Stwierdzenie \ref{prop:kernel} ma 2 skutki praktyczne. Kolejność stosowania reguł PK nie
zmienia wyniku. Sprawdzenie tekstu rozkłada się na niezależne kontrole: 1 kontrola
dla 1 reguły.

Cel standardu ma w modelu 2 warunki. Warunek spójności: $\sem{s}'\neq\emptyset$ dla
każdego $s\in L'$, czyli konwencje zostawiają zdaniu co najmniej 1 odczytanie. Warunek
jednoznaczności: $|\sem{s}'|\le 1$ dla każdego $s\in L'$. Przy obu warunkach relacja $R'$
jest funkcją.

\begin{remark}[status celu]
\label{rem:gap}
Standard PK 0.1 formułuje cel słowami: czytelnik, który zna konwencję, odczytuje każde
zdanie w 1 sposób. Cel jest hipotezą standardu. Numerowane reguły PK 0.1 celu nie
gwarantują. Zdanie \pkq{Operatorowi NIE WOLNO włączyć pompy P1 i pompy P2.} spełnia każdą
regułę numerowaną. Zdanie ma 2 odczytania: zakaz dotyczy każdej z 2 pomp albo zakaz dotyczy
włączenia 2 pomp naraz. Zdanie \pkq{Każdy operator obsługuje dokładnie 1 pompę.} także
spełnia reguły i ma 2 odczytania zasięgu. Standard zamyka takie przypadki krokiem 7
kontroli, czyli testem 2 znaczeń. Test 2 znaczeń wymaga oceny znaczenia przez człowieka.
Pozycje G1 i G2 zbioru testowego mierzą te 2 luki (sekcja \ref{sec:probe}).
\end{remark}

\subsection{Miejsca wyboru}

\begin{definition}
\emph{Miejsce wyboru} zdania $s$ to konstrukcja, która dopuszcza co najmniej 2 odczytania
przy ustalonych odczytaniach reszty zdania. Miejsce wyboru $i$ ma zbiór wariantów $C_i$.
Miejsca wyboru są \emph{niezależne}, gdy każda kombinacja wariantów daje inne odczytanie.
\end{definition}

\begin{theorem}
\label{thm:product}
Jeżeli zdanie $s$ ma $k$ niezależnych miejsc wyboru, to
$a(s)=\prod_{i=1}^{k}|C_i|$ oraz $\log_2 a(s)=\sum_{i=1}^{k}\log_2|C_i|$.
Zamknięcie miejsca wyboru $i$ zmniejsza $\log_2 a(s)$ o $\log_2|C_i|$ bitów.
\end{theorem}

\begin{proof}
Zbiór odczytań jest w bijekcji z iloczynem kartezjańskim $C_1\times\dots\times C_k$.
\end{proof}

Twierdzenie \ref{thm:product} tłumaczy, dlaczego reguły PK działają osobno. Każda reguła
zamyka 1 typ miejsca wyboru i odejmuje stałą liczbę bitów. Miejsca wyboru bywają zależne;
wtedy liczba odczytań jest mniejsza niż iloczyn. Pomiar P1 liczy odczytania dokładnie dla
4 konstrukcji (tabela \ref{tab:counts}, rysunek \ref{fig:growth}).

Według reguły B8 zdanie opisu ma co najwyżej 25 słów. Samo ograniczenie długości nie usuwa
niejednoznaczności. Liczba nawiasowań binarnych ciągu 25 słów to liczba Catalana
$C_{24}=\text{\pkCatalanTwentyFour}$. Zbiór odczytań jest więc skończony, a górna granica liczby odczytań
ma rząd $10^{12}$. Jednoznaczność dają dopiero pozostałe reguły.

\subsection{Strefa szara}

\begin{definition}
\emph{Strefa szara} zdania $s$ to zbiór
$G(s)=\bigcup\sem{s}\setminus\bigcap\sem{s}$.
\end{definition}

Strefa szara zawiera sytuacje, w których odczytania dają różną wartość logiczną. Poza
strefą szarą odczytania są zgodne. Strefa szara jest pusta dokładnie wtedy, gdy każde
2 odczytania zdania są równe; dla zdania jednoznacznego $G(s)=\emptyset$.

\begin{example}
Zdanie \pkq{Dzieci do 3 lat wchodzą bezpłatnie.} ma odczytania $m_1=\{w:\text{wiek}<3\}$
i $m_2=\{w:\text{wiek}<4\}$. Strefa szara to $\{w:3\le\text{wiek}<4\}$. Dziecko, które ma
3 lata i 5 miesięcy, leży w strefie szarej.
\end{example}

Model strefy szarej odpowiada nadwartościowaniu Fine'a \cite{fine1975}: zdanie jest
nadprawdziwe w $\bigcap\sem{s}$ i nadfałszywe poza $\bigcup\sem{s}$. Rysunek
\ref{fig:gray} pokazuje strefy szare 5 pozycji testowych, a rysunek \ref{fig:truth}
pokazuje ten sam obiekt w tablicach prawdy. Pomiar P4 opiera się na strefie szarej: sonda
pyta model o sytuację, w której odczytania się różnią.

\rysunek{gray_zone}{Semantyka: 2 odczytania jako zbiory na osi liczbowej i strefa szara,
w której odczytania się różnią.}{fig:gray}

\rysunek{truth_tables}{Logika: tablice prawdy 2 odczytań dla reguł C2, C1 i C4.}{fig:truth}

\subsection{Moc zbioru odczytań}

Zbiór odczytań ma w modelu 3 możliwe moce.
\begin{itemize}[leftmargin=1.6em]
  \item \textbf{Moc 1.} Zdanie jest jednoznaczne. Taki stan jest celem PK.
  \item \textbf{Moc skończona większa niż 1.} Źródłem jest składnia albo słownik: zasięg
  określenia, zaimek, spójnik, granica przedziału. Lista zakończona skrótem \pkq{itp.}
  oznacza nadzbiór zbioru wymienionych elementów; zdanie nie wskazuje, który nadzbiór. Dla dziedziny o $u$ elementach i listy
  o $k$ elementach zdanie ma $2^{u-k}$ odczytań; dla $u=10$ i $k=2$ zdanie ma \pkOpenList{} odczytań.
  \item \textbf{Moc continuum.} Słowo nieostre, na przykład \pkq{regularnie}, dopuszcza
  każdy próg z przedziału liczb rzeczywistych. Zbiór dopuszczalnych uściśleń jest
  nieprzeliczalny \cite{fine1975}. Reguła A5 każe wybrać 1 próg i zapisać próg liczbą
  z jednostką.
\end{itemize}

Osobną klasę tworzy deiksa. Znaczenie słowa \pkq{jutro} jest funkcją kontekstu; Kaplan
nazywa taką funkcję charakterem wyrażenia \cite{kaplan1989}. Reguła D3 wymaga daty.
Zdanie z datą ma stały charakter: znaczenie zdania jest takie samo w każdym dniu czytania.

\subsection{Klasyfikacja reguł}
\label{sec:rules}

Autor przypisał każdej z \pkRules{} reguł 3 cechy: rodzaj operatora, klasę sprawdzarki
i mechanizm redukcji kosztu. Klasyfikację zawiera plik \texttt{data/rules.csv};
rysunek \ref{fig:rules} pokazuje wynik. Ograniczeniem jest \pkRulesR{} reguł.
Ograniczenie i konwencję łączy \pkRulesRK{} reguł. Samą konwencją jest \pkRulesK{} reguła:
reguła E6, która każe ujawnić konwencje w legendzie. Klasyfikacja jest oceną autora
i nie ma drugiego oceniającego.

\rysunek{rules}{Klasyfikacja reguł PK 0.1 według rodzaju operatora, klasy sprawdzarki
i mechanizmu redukcji kosztu.}{fig:rules}

\section{Złożoność obliczeniowa}
\label{sec:complexity}

\subsection{Zadania czytelnika i mechanizmy PK}

Czytelnik tekstu wykonuje 4 zadania. Zadanie PARSUJ wyznacza zbiór odczytań zdania.
Zadanie WYBIERZ wskazuje odczytanie, które zamierzył autor. Zadanie SPRAWDŹ rozstrzyga,
czy tekst spełnia reguły. Zadanie WNIOSKUJ rozstrzyga, czy zbiór zdań jest spójny i co ze
zbioru wynika. PK zmienia koszt tych zadań 3 mechanizmami.
\begin{itemize}[leftmargin=1.6em]
  \item \textbf{M1, zwinięcie zbioru odczytań.} Reguła zmniejsza liczbę $a(s)$ do wartości 1. Zadanie
  WYBIERZ znika.
  \item \textbf{M2, ograniczenie rozmiaru.} Reguły B8, B4, C2 i C5 ograniczają liczbę słów,
  liczbę rzeczowników w ciągu, mieszanie spójników i liczbę przeczeń. Koszt zadania
  PARSUJ dla 1 zdania ma stałą górną granicę, więc koszt dla tekstu rośnie liniowo
  z liczbą zdań.
  \item \textbf{M3, lokalność.} Reguły A2, A3, D3 i E4 usuwają zaimek, deiksę i odesłanie
  bez numeru. Czytelnik interpretuje zdanie bez przeszukiwania reszty tekstu. Tekst jest
  koniunkcją zdań ze wspólnym słownikiem nazw.
\end{itemize}
Kosztem PK jest długość: nazwa zastępuje zaimek, liczba z jednostką zastępuje słowo
nieostre, a 2 zdania zastępują imiesłów.

\subsection{Zadanie PARSUJ}
\label{sec:parse}

Dla gramatyki bezkontekstowej algorytm Earleya działa w czasie rzędu $n^3$; dla gramatyki
jednoznacznej w czasie rzędu $n^2$ \cite{earley1970}. Dla gramatyki LR($k$) istnieje parser
liniowy \cite{knuth1965,leo1991}. Swobodny szyk podnosi koszt: problem rozpoznawania dla
gramatyk ID/LP jest NP-zupełny \cite{barton1985}. Reguła B2 ustala szyk: podmiot,
orzeczenie, dopełnienie. Reguła B6 stawia określenie przy słowie określanym.

Pomiar P2 porównuje 2 gramatyki modelowe. Gramatyka pierwsza pozwala przyłączyć frazę
przyimkową do czasownika albo do rzeczownika, który stoi wcześniej. Gramatyka druga,
zgodna z regułą B6, przyłącza frazę do rzeczownika bezpośrednio przed frazą. Obie gramatyki
generują te same ciągi słów. Wynik podaje sekcja \ref{sec:resparse}. Gramatyki są modelami
1 konstrukcji i nie opisują całej polszczyzny.

\subsection{Zadanie WYBIERZ}

\begin{theorem}
\label{thm:questions}
Niech $Z$ oznacza zamierzone odczytanie zdania $s$, czyli zmienną losową o rozkładzie $P$
na zbiorze $\sem{s}$. Każda strategia pytań z odpowiedzią \pkq{Tak} albo \pkq{Nie}, która zawsze
wskazuje $Z$, zadaje średnio co najmniej $H(Z)=-\sum_m P(m)\log_2P(m)$ pytań. Istnieje
strategia, która zadaje średnio mniej niż $H(Z)+1$ pytań.
\end{theorem}

\begin{proof}
Strategia pytań jest drzewem binarnym, czyli kodem prefiksowym dla zbioru $\sem{s}$.
Oba oszacowania wynikają z twierdzenia Shannona o kodowaniu bez szumu \cite{shannon1948}.
\end{proof}

Wielkość $H(Z)$ jest brakiem informacji w zdaniu. Czytelnik pobiera brakujące bity
z kontekstu, z pytania do autora albo z własnego domysłu. Dla zdania PK brak wynosi 0 bitów.
Dla zdania z 2 równie prawdopodobnymi odczytaniami brak wynosi 1 bit. Czytelnik, który
zgaduje, myli się wtedy w 50 \% przypadków przy każdej strategii zgadywania.

\subsection{Zadanie WNIOSKUJ}
\label{sec:horn}

Praca definiuje rdzeń zdaniowy PK dla wymagań. Rdzeń obejmuje wzór \pkq{Jeżeli W, to D.}
z reguły C3, zakaz mieszania spójników z reguły C2 i limit przeczeń z reguły C5.

\begin{definition}
Niech $A$ oznacza zbiór zdań prostych twierdzących, czyli \emph{atomów}. \emph{Literał} to
atom albo atom z przeczeniem. Język \PKz{} zawiera zdania 2 postaci:
\begin{center}
\pkq{Jeżeli $l_1$ i \dots{} i $l_k$, to $a$.}\qquad
\pkq{Jeżeli $l_1$ i \dots{} i $l_k$, to konflikt.}
\end{center}
gdzie $k\ge0$, każde $l_i$ jest literałem, a $a$ jest atomem. Zdanie drugiej postaci
zapisuje 2 wymagania dla 1 działania: \pkq{MUSI} oraz \pkq{NIE WOLNO}. Język \PKzj{j}
zawiera zdania, w których poprzednik ma co najwyżej $j$ przeczeń. Zbiór zdań $\Gamma$ jest
\emph{spójny}, gdy co najmniej 1 wartościowanie atomów spełnia każde zdanie zbioru $\Gamma$;
zdanie drugiej postaci jest spełnione, gdy poprzednik jest fałszywy.
\end{definition}

\begin{theorem}
\label{thm:horn}
Spójność zbioru zdań języka \PKzj{0} jest rozstrzygalna w czasie liniowym względem
łącznej długości zdań.
\end{theorem}

\begin{proof}
Zdanie \pkq{Jeżeli $a_1$ i \dots{} i $a_k$, to $a$.} jest klauzulą
$\neg a_1\lor\dots\lor\neg a_k\lor a$. Zdanie z konfliktem jest klauzulą
$\neg a_1\lor\dots\lor\neg a_k$. Każda klauzula ma co najwyżej 1 literał dodatni, czyli
jest klauzulą Horna. Algorytm Dowlinga i Galliera rozstrzyga spełnialność formuły Horna
w czasie liniowym \cite{dowling1984}.
\end{proof}

\begin{theorem}
\label{thm:np}
Spójność zbioru zdań języka \PKzj{1} jest problemem NP-zupełnym.
\end{theorem}

\begin{proof}
Problem należy do klasy NP: świadkiem jest wartościowanie. NP-trudność wynika z redukcji
problemu 3-SAT \cite{cook1971}. Niech $C$ będzie klauzulą o 3 literałach nad atomami
$x$, $y$, $z$. Redukcja zależy od liczby literałów dodatnich w klauzuli $C$.
\begin{center}
\begin{tabular}{ll}
\toprule
Klauzula $C$ & Zdania języka \PKzj{1}\\
\midrule
$\neg x\lor\neg y\lor\neg z$ & Jeżeli $x$ i $y$ i $z$, to konflikt.\\
$\neg x\lor\neg y\lor z$ & Jeżeli $x$ i $y$, to $z$.\\
$\neg x\lor y\lor z$ & Jeżeli $x$ i nie $y$, to $z$.\\
$x\lor y\lor z$ & Jeżeli nie $x$, to $u$. Jeżeli $u$ i nie $y$, to $z$.\\
\bottomrule
\end{tabular}
\end{center}
W 3 pierwszych wierszach zdanie jest równoważne klauzuli. W wierszu czwartym $u$ jest
nowym atomem, a zdania dają klauzule $x\lor u$ oraz $\neg u\lor y\lor z$. Jeżeli
wartościowanie spełnia $x\lor y\lor z$, to rozszerzenie $u=\neg x$ spełnia obie klauzule.
W drugą stronę rezolucja względem $u$ daje $x\lor y\lor z$. Każde zdanie ma co najwyżej
1 przeczenie, a liczba zdań jest co najwyżej 2 razy większa niż liczba klauzul.
\end{proof}

Skrypt sprawdził redukcję na \pkRedInstances{} losowych formułach, z których każda ma
co najwyżej 6 zmiennych: pełne wyliczenie wartościowań dało zgodną odpowiedź w każdym
przypadku. Twierdzenie \ref{thm:np} zgadza się z dychotomią Schaefera \cite{schaefer1978}:
zbiór klauzul Horna z klauzulą $x\lor u$ leży poza każdą klasą łatwą.

\begin{remark}[znaczenie dla standardu]
Reguła C5 dopuszcza 1 przeczenie w zdaniu, a reguła C3 zaleca wzór \pkq{Jeżeli nie W,
to D2.} dla przypadku przeciwnego. Zbiór wymagań zgodny z PK 0.1 jest więc zbiorem zdań
języka \PKzj{1}. Zbiór bez przeczeń w poprzednikach leży po stronie czasu liniowego.
Zbiór z 1 przeczeniem w poprzedniku leży po stronie NP-zupełności. NP-zupełność opisuje przypadek najgorszy; programy
rozstrzygające spełnialność radzą sobie z praktycznymi zbiorami wymagań. Pełna PK 0.1
nie ogranicza ekspresywności polszczyzny, więc wnioskowanie w pełnej PK jest
nierozstrzygalne tak samo jak w logice pierwszego rzędu. Rysunek \ref{fig:classes}
zestawia wyniki z fragmentami języka angielskiego \cite{pratt2006}.
\end{remark}

\rysunek{classes}{Teoria złożoności: położenie problemów związanych z PK wśród klas
złożoności.}{fig:classes}

\subsection{Zadanie SPRAWDŹ}
\label{sec:check}

Reguły PK różnią się kosztem kontroli. Autor wyróżnił 3 klasy sprawdzarek
(rysunek \ref{fig:rules}). Automat skończony nad ciągiem słów rozstrzyga \pkRulesRegular{}
z \pkRules{} reguł: listę słów zakazanych, długość zdania, liczbę przeczeń, zapis granicy.
Analizy składni wymaga \pkRulesSyntactic{} reguł, na przykład szyk z reguły B2 i zasięg
z reguły B6. Oceny znaczenia wymaga \pkRulesSemantic{} reguł, na przykład reguła A1:
1 termin dla 1 pojęcia.

Plik \texttt{src/pklint.py} zawiera sprawdzarkę z \pkLintChecks{} kontrolami. Każda
kontrola jest wyrażeniem regularnym nad słowami 1 zdania, więc sprawdzarka działa
w czasie liniowym. Sprawdzarka wskazuje kandydatów; ocenę kończy człowiek. Na zbiorze
testowym sprawdzarka wykrywa \pkLintDetected{} z \pkItems{} zdań zwykłych i nie zgłasza
żadnego z \pkItemsPkSentences{} zdań PK. Pozostałe \pkLintUndetected{} zdań zwykłych wymaga analizy składni
albo znaczenia.

\subsection{Koszt w modelu językowym}
\label{sec:cost}

Koszt przejścia tekstu przez transformator zależy od liczby tokenów $n$: warstwa uwagi
wykonuje pracę rzędu $n^2$ \cite{vaswani2017}. Rozstrzygnięcie niejednoznaczności tej pracy
nie zmienia. Merrill i Sabharwal wykazali, że transformator o precyzji logarytmicznej
liczy tylko funkcje z klasy TC$^0$ \cite{merrill2023}. Model przenosi więc trudniejsze
obliczenia do kolejnych generowanych tokenów. Zdanie, które wymaga wyboru odczytania,
dokłada modelowi pracy; zdanie jednoznaczne tej pracy nie dokłada. Praca nie mierzy tego
efektu wewnątrz modelu; pomiar P4 mierzy zachowanie modelu na wyjściu.

\begin{theorem}
\label{thm:breakeven}
Niech tekst zwykły ma $n$ tokenów, a tekst PK ma $(1+\delta)n$ tokenów. Niech czytelnik
odczytuje tekst zwykły błędnie z prawdopodobieństwem $p$, a tekst PK z prawdopodobieństwem
$e$. Niech po błędnym odczytaniu zadanie wraca do autora i czytelnik czyta tekst od nowa,
a próby są niezależne. Oczekiwany koszt odczytania w tokenach wynosi $n/(1-p)$ dla tekstu
zwykłego i $(1+\delta)n/(1-e)$ dla tekstu PK. Tekst PK jest tańszy dokładnie wtedy, gdy
\begin{equation}
\label{eq:breakeven}
p>1-\frac{1-e}{1+\delta}.
\end{equation}
\end{theorem}

\begin{proof}
Liczba prób ma rozkład geometryczny o wartości oczekiwanej $1/(1-p)$ albo $1/(1-e)$.
Nierówność \eqref{eq:breakeven} wynika z porównania 2 kosztów.
\end{proof}

Dla $e=0$ warunek \eqref{eq:breakeven} ma postać $p>\delta/(1+\delta)$. Model pomija koszt
wykrycia błędu i koszt szkody; oba koszty przesuwają próg na korzyść PK. Model pomija
także pracę autora. Autor pisze tekst 1 raz, a czytelnicy czytają tekst $k$ razy, więc
dodatkowa praca autora rozkłada się na $k$ odczytów. Model językowy czyta instrukcję
systemową przy każdym wywołaniu.
